\chapter{Progettazione e sviluppo del prototipo}
\label{chap:progettazione-prototipo}

Questo capitolo copre la seconda fase dello stage: a partire dagli otto criteri \emph{\emph{core}} individuati nel capitolo~\ref{chap:metodologia-selezione}, vengono qui presentati i requisiti dell'estensione, le scelte di progettazione e architettura, la strategia di prompting adottata per integrare un \acr{llm} nella verifica, e le fasi di prompt engineering, sperimentazione e confronto tra modelli che hanno portato al consolidamento del prototipo, in preparazione della validazione sperimentale discussa nel capitolo~\ref{chap:validazione-benchmark}.

\section{Analisi dei requisiti}
\label{sec:analisi-requisiti}

\todoplaceholder{requisiti funzionali e non funzionali, casi d'uso, analisi dei rischi.}

% TODO: requisiti funzionali e non funzionali dell'estensione, casi d'uso
% principali, eventuale analisi preventiva dei rischi. Richiamare qui, per i
% soli otto criteri \emph{core}, la classificazione di dettaglio già presentata nel
% capitolo~\ref{chap:metodologia-selezione}.


\section{Architettura e scelte di progettazione}
\label{sec:progettazione}

\todoplaceholder{architettura dell'estensione, scelte tecnologiche, integrazione con l'LLM.}

% TODO: architettura dell'estensione browser, scelte tecnologiche, modalità di
% integrazione con il modello linguistico, design pattern adottati.


\section{Strategia di prompting e primo prototipo}
\label{sec:strategia-prompting}

\todoplaceholder{definizione della strategia di prompting, utilizzo di un frontier LLM, sviluppo dei primi controlli, primo prototipo, primi test.}

% TODO (settimana 2 del piano di lavoro): definizione della strategia di
% prompting per i controlli assistiti da LLM sui criteri \emph{core}; scelta e
% utilizzo di un frontier LLM come riferimento iniziale; sviluppo dei primi
% controlli automatici/assistiti; realizzazione del primo prototipo
% funzionante; esito dei primi test esplorativi.


\section{Prompt engineering, sperimentazione e confronto tra modelli}
\label{sec:prompt-engineering-confronto}

\todoplaceholder{ottimizzazione dei prompt, sperimentazione con modelli LLM più piccoli, confronto tra modelli, consolidamento del prototipo.}

% TODO (settimana 3 del piano di lavoro): attività di prompt engineering e
% ottimizzazione iterativa dei prompt definiti nella sezione precedente;
% sperimentazione con modelli LLM di dimensioni minori rispetto al frontier
% model iniziale; confronto qualitativo/quantitativo fra i modelli provati
% (obiettivo desiderabile); consolidamento del prototipo risultante.


\section{Implementazione}
\label{sec:implementazione}

\todoplaceholder{aspetti realizzativi più significativi, con frammenti di codice.}

% TODO: aspetti realizzativi più significativi, con frammenti di codice
% commentati dove utile.


\section{Consolidamento del prototipo e preparazione della validazione}
\label{sec:consolidamento-prototipo}

\todoplaceholder{stato finale del prototipo al termine della fase di sviluppo, e preparazione della fase di validazione sperimentale.}

% TODO: sintesi dello stato del prototipo consolidato al termine della
% settimana 3, come base per la validazione sperimentale e la costituzione
% del benchmark descritte nel capitolo~\ref{chap:validazione-benchmark}.